% SPDX-License-Identifier: Apache-2.0
\section{Validation plan}
\label{sec:validation}

The plan below is not aspirational. Every step of the certification pipeline --- split, band,
calibration, certificate, held-out validation --- was executed in this study on public data, and
the results are in Section~\ref{sec:results}; what a PoC sprint changes is the data, not the
protocol. Some pre-registered items were \emph{not} executed, and are named here rather than
left to be discovered: the two kernel hypotheses and the entanglement ablation, all closed by
the screens before they could run, and the drift measurements recorded in amendment~A7, for
which no script was written. The appendix enumerates both. A validation plan that cannot say
which of its own steps ran is not a plan, and it is why the negative quantum results are in the
body rather than buried.

\textbf{The success metric is a certificate, not a score.} A PoC succeeds if the
band-conditional false-decline rate is certified at a level the business names, with
family-wise error control, and the certificate then holds on a block that was never used to
build it. Average precision and ROC AUC are reported alongside, but they decide nothing: they
are not statements about a decision and they carry no finite-sample guarantee. The
corresponding failure is equally well defined --- the admissible set comes back empty, which is
a real outcome and is what happens here at the tightest band budget.

\textbf{Held-out validation is the acceptance test.} The certificate is selected on
calibration data and then read once on a fold protected by a persisted ledger that refuses a
second distinct configuration. In this study \ClaimHFiveConfigsHolding{} of
\ClaimHFiveConfigsTotal{} certified configurations held their risk on that fold, the worst
realised risk being \ClaimHFiveWorstRatio{} of its budget. A PoC would run exactly this test
and report it whichever way it came out.

\textbf{Benchmarking is against a tuned baseline, replicated, with the power stated first.}
The comparator is a tuned gradient-boosted model, not a straw man, fitted on identical rows and
features. Intervals resample entities rather than rows, because transactions on one card are
not independent. Crucially, a detectable-effect ceiling is fixed \emph{before} the comparison: here
\ClaimPowerCeiling{}, against an effect the comparison could actually resolve of
\ClaimMpsMdeRealised{}, computed afterwards from its own bootstrap standard error. That gap is
why H4 is reported as underpowered rather than as a null. A PoC would size its evaluation block to that ceiling
first, which is the single change that would make a quantum-versus-classical comparison on this
problem informative.

\textbf{What keeps the plan honest.} Hypotheses, grids, splits and decision thresholds are
frozen and hashed before any model is fitted, and a gate refuses undocumented change; multiple
comparisons carry Holm correction; coverage is judged against the exact law of
Section~\ref{sec:method} rather than against $\alpha$. Every retraction is recorded --- this
submission carries \ClaimDecisionEntries{} decision entries, including the ones that removed
results we had already written down --- and reproduction is
\texttt{make smoke \&\& make reproduce \&\& make check}.
